% !TeX root = ../main.tex

\section{Literature Review} \label{sec: literature review}
This section will review the existing literature on the topic and work out the research gap this thesis aims to explore. Section \ref{subsec: literature_commodity forecasting} goes over existing approaches used in commodity forecasting and presents previous and remaining challenges. Section \ref{subsec: literature_transfer learning} explores the field of transfer learning as a potential solution for some of these remaining challenges. Section \ref{subsec: literature_research gap} shows where the existing literature leaves a gap that this thesis aims to fill, as the intersection of transfer learning, commodity markets and data-sparse emerging economies.

%-----------------------------------------------------------------------------------------
\subsection{Commodity Forecasting}  \label{subsec: literature_commodity forecasting}
Before evaluating whether transfer learning improves commodity forecasting in data-sparse markets, it is necessary to establish what classical, single-market methods already achieve when data is not scarce. This benchmark matters for two reasons: classical models supply the target-only baselines agains which the proposed framework is compared, see Section \ref{subsec: methodology_baseline models}, \orangemarker{and their known limitations, even under the most favourable data conditions, motivate the case for knowledge transfer in the first place}.

\subsubsection{Classical Approaches}    \label{subsubsec: literature_classical approaches}
A large body of empirical literature deals with regressions of commodity series on macroeconomic and financial predictors and evaluates the in-sample and out-of-sample forecasting performance. For instance, \textcite{garganoForecastingCommodityPrice2014} study the predictive power over different forecast horizons, economic states, and types of commodities. Using Reuters/Jefferies-CRB commodity indexes over 1947-2010, they find that predictability from variables such as industrial production growth, commodity-currency exchange rates and global economic activity indices is present but varies highly between different commodities. While it is strongest for industrials, metals and the broad commodity index, and weaker for fats/oils, foods and livestock. Furthermore, they find forecasting performance to be more pronounced during recessions. Previously, \textcite{chenCanExchangeRates2010} showed that the exchange rates of only five commodity exporters (Australia, Canada, New Zealand, South Africa and Chile) have solid forecasting power over a global aggregate commodity price index made up by more than forty products traded on various exchanges and published by the IMF. Further studies find more limited results for the predictability of commodity spot indices through commodity exchange rates and other macroeconomic variables \parencite{groenCommodityPricesCommodity2011}. \textcite{guidolin_forecasting_2021} examine whether commodity-specific factors (basis, hedging pressure and momentum) improve forecasting performance. The idea is to test whether they include information not already captured by common macroeconomic variables. While including commodity-specific variables does not improve the forecasting perfomance, their importance seems to depend on whether statistical or economic loss functions are used. This apparent disconnect between statistical forecast accuracy of returns and the economic value created when included in a portfolio choice problem is also present in other studies \parencite{cenesizogluReturnPredictionModels2012, guidolinPredictingCommodityReturns2026}. 

\subsubsection{Machine Learning \& Deep Learning in Commodity Markets}
Over the past two decades commodity forecasting has shifted from classical statistical approaches toward machine learning and deep learning techniques. Traditional models such as \ac{arima}, \ac{garch} and exponential smoothing rely on assumptions of stationarity and linearity that frequently fail for commodity price series characterized by nonlinearity, volatility clustering and structural breaks \parencite{foroutanDeepLearningSystems2024,sunNovelAgriculturalCommodity2025}. Machine learning methods like \ac{svr} or \ac{xgboost} provide greater flexibility in modeling multidimensional relationships without imposing strict distributional assumptions \parencite{manognaNovelHybridNeural2025}. Deep learning architectures, particularly \ac{mlp}, \ac{rnn}, \ac{lstm}, \ac{gru}, \ac{esn} and \ac{bilstm} have further advanced the field by capturing long-term temporal dependencies and complex nonlinear patterns \parencite{manognaNovelHybridNeural2025,sunNovelAgriculturalCommodity2025}. Next to individual approaches, hybrid models that combine decomposition techniques, ensemble methods and multiple neural architectures have emerged as a dominant trend. \textcite{fangOptimalForecastCombination2020} use a combination of three models, \ac{svm}, \ac{nn} and \ac{arima}, to produce short-term forecasts and find that it outperforms the individual ones. \textcite{rayARIMALSTMModelPredicting2023} report similar results with a comparable ensemble model using a monthly forecasting period.

The literature spans multiple commodity sectors, with agricultural prices receiving a lot of attention due to their volatility and implications for food security \parencite{manognaNovelHybridNeural2025,sariVariousOptimizedMachine2024}. Energy commodities, especially crude oil, constitute a second cluster, driven by their geopolitical importance and market liquidity \parencite{foroutanDeepLearningSystems2024,senForecastingCrudeOil2024,liangCrudeOilPrice2023}. Studies like those from \textcite{varshiniHowGoodAre2024} or \textcite{jinMachineLearningGold2025} cover forecasting of precious and industrial metals as another category, with both general economic and financial importance. 

\subsection{Data Scarcity in Emerging Economies}
\subsection{Cross-Market Commodity Dynamics}


%-----------------------------------------------------------------------------------------
\subsection{Transfer Learning: Foundations \& Applications} \label{subsec: literature_transfer learning}
This shortcoming naturally invites the exploration of methodologies, which aim to built predictors, that allow for differences cross-market. The goal is to leverage data-rich markets, while still being able to include market-specific characteristics. This approach can be categorised as a transfer learning technique, which is a subset of machine learning, that allows for the assumption of different domains, tasks and distributions between the training and test set \parencite{panSurveyTransferLearning2010}.
\subsubsection{What is Transfer Learning?}
\subsubsection{Transfer Learning in Finance}
\subsubsection{Domain Adaptation Techniques}
\textcite{mashettyTransferLearningCrossMarket2025a} \orangemarker{explain approach} This approach proved to outperform both a traditional machine learning model and a deep learning model, which were trained only on target data, in all three chosen metrics, namely \ac{mae}, \ac{rmse} and directional accuracy.

%-----------------------------------------------------------------------------------------
\subsection{Research Gap \& Positioning}    \label{subsec: literature_research gap}
